\documentclass[11pt]{article}
\usepackage[margin=1in]{geometry}
\usepackage{amsmath}
\usepackage{amssymb}
\usepackage{amsthm}
\usepackage{enumitem}
\usepackage{xcolor}
\usepackage{tcolorbox}

\title{Detailed Explanation of Question 1c\\ORF 309 S2025 Midterm 1}
\author{}
\date{}

\begin{document}

\maketitle

\section{Problem Setup}

Bob takes an exam with 2 questions, each worth 10 points.

\textbf{Given information:}
\begin{itemize}
    \item $P(A) = 0.8$ (prob of getting question 1 correct)
    \item $P(B|A) = 0.6$ (prob of getting question 2 correct given question 1 correct)
    \item $P(B|A^c) = 0.1$ (prob of getting question 2 correct given question 1 wrong)
\end{itemize}

\textbf{From parts (a) and (b):}
\begin{itemize}
    \item $P(B) = 0.5$ (total prob of getting question 2 correct)
    \item $P(A \cap B) = P(B|A)P(A) = 0.6 \times 0.8 = 0.48$
\end{itemize}

\section{Part (c): Find Variance of Bob's Grade}

\subsection{Define the Random Variable}

Let $X$ be Bob's total grade. Using indicator random variables:
\begin{equation}
X = 10 \cdot (1_A + 1_B)
\end{equation}

where:
\begin{itemize}
    \item $1_A = \begin{cases} 1 & \text{if Bob gets question 1 correct} \\ 0 & \text{otherwise} \end{cases}$
    \item $1_B = \begin{cases} 1 & \text{if Bob gets question 2 correct} \\ 0 & \text{otherwise} \end{cases}$
\end{itemize}

\subsection{Step 1: Calculate $E[X]$}

By linearity of expectation:
\begin{align}
E[X] &= 10 \cdot E[1_A + 1_B]\\
&= 10 \cdot (E[1_A] + E[1_B])\\
&= 10 \cdot (P(A) + P(B))\\
&= 10 \times (0.8 + 0.5)\\
&= 13
\end{align}

\textbf{Key fact:} For any indicator random variable $1_A$, we have $E[1_A] = P(A)$.

\subsection{Step 2: Calculate $E[X^2]$ --- The Critical Step}

This is where most students struggle. Let's do it step by step.

\begin{equation}
E[X^2] = E\left[(10(1_A + 1_B))^2\right] = 10^2 \cdot E\left[(1_A + 1_B)^2\right]
\end{equation}

\subsubsection{Expand the Square}

\begin{tcolorbox}[colback=yellow!10!white,colframe=orange!75!black,title=Key Algebraic Step]
\begin{equation}
(1_A + 1_B)^2 = 1_A^2 + 1_B^2 + 2 \cdot 1_A \cdot 1_B
\end{equation}
\end{tcolorbox}

\subsubsection{Simplify Using Indicator Properties}

For any indicator random variable $1_A$:
\begin{tcolorbox}[colback=blue!5!white,colframe=blue!75!black,title=Important Property]
\begin{equation}
1_A^2 = 1_A
\end{equation}

\textbf{Why?} Because $1_A \in \{0, 1\}$, so $0^2 = 0$ and $1^2 = 1$.
\end{tcolorbox}

Therefore:
\begin{equation}
(1_A + 1_B)^2 = 1_A + 1_B + 2 \cdot 1_A \cdot 1_B
\end{equation}

\subsubsection{Take Expectation}

\begin{align}
E[(1_A + 1_B)^2] &= E[1_A + 1_B + 2 \cdot 1_A \cdot 1_B]\\
&= E[1_A] + E[1_B] + 2E[1_A \cdot 1_B]
\end{align}

\textbf{Key question:} What is $E[1_A \cdot 1_B]$?

\begin{tcolorbox}[colback=green!5!white,colframe=green!75!black,title=Product of Indicators]
For indicator random variables:
\begin{equation}
1_A \cdot 1_B = 1_{A \cap B}
\end{equation}

\textbf{Why?} The product equals 1 only when \textbf{both} $A$ and $B$ occur.

Therefore:
\begin{equation}
E[1_A \cdot 1_B] = E[1_{A \cap B}] = P(A \cap B)
\end{equation}
\end{tcolorbox}

\subsubsection{Substitute Values}

\begin{align}
E[(1_A + 1_B)^2] &= P(A) + P(B) + 2P(A \cap B)\\
&= 0.8 + 0.5 + 2 \times 0.48\\
&= 1.3 + 0.96\\
&= 2.26
\end{align}

\subsubsection{Final Calculation of $E[X^2]$}

\begin{equation}
E[X^2] = 10^2 \cdot E[(1_A + 1_B)^2] = 100 \times 2.26 = 226
\end{equation}

\subsection{Step 3: Calculate Variance}

Using the variance formula:
\begin{tcolorbox}[colback=red!5!white,colframe=red!75!black,title=Variance Formula]
\begin{equation}
\text{Var}(X) = E[X^2] - (E[X])^2
\end{equation}
\end{tcolorbox}

\begin{align}
\text{Var}(X) &= 226 - 13^2\\
&= 226 - 169\\
&= 57
\end{align}

\textbf{Note:} The solution PDF has a typo showing $13^2 = 132$, but the correct value is $13^2 = 169$. The final answer of 57 is still correct.

\section{Alternative Method: Using Variance of Sum Formula}

This method uses the general formula from section 3.2 and is often more intuitive than computing $E[X^2]$ directly.

\subsection{Reframe the Problem}

Since $X = 10 \cdot (1_A + 1_B)$, we can define:
\begin{itemize}
    \item $X_1 = 10 \cdot 1_A$ (points from question 1)
    \item $X_2 = 10 \cdot 1_B$ (points from question 2)
    \item Then $X = X_1 + X_2$
\end{itemize}

\subsection{Step 1: Find $\text{Var}(X_1)$ and $\text{Var}(X_2)$}

For an indicator variable $1_A$, the variance is:
\begin{equation}
\text{Var}(1_A) = P(A)(1 - P(A)) = 0.8 \times 0.2 = 0.16
\end{equation}

\begin{tcolorbox}[colback=cyan!5!white,colframe=cyan!75!black,title={Why Does $\text{Var}(1_A) = P(A)(1-P(A))$?}]
For any indicator variable $1_A$ that takes value 1 with probability $P(A)$ and 0 with probability $1-P(A)$:

\textbf{Step 1: Calculate $E[1_A]$}
\begin{equation*}
E[1_A] = 1 \cdot P(A) + 0 \cdot P(A^c) = P(A)
\end{equation*}

\textbf{Step 2: Calculate $E[1_A^2]$}
\begin{equation*}
E[1_A^2] = E[1_A] = P(A) \quad \text{(since $1_A^2 = 1_A$ for indicators)}
\end{equation*}

\textbf{Step 3: Apply variance formula}
\begin{align*}
\text{Var}(1_A) &= E[1_A^2] - (E[1_A])^2\\
&= P(A) - (P(A))^2\\
&= P(A) - P(A)^2\\
&= P(A)(1 - P(A))
\end{align*}

This is the standard formula for the variance of a Bernoulli random variable.
\end{tcolorbox}

Similarly:
\begin{equation}
\text{Var}(1_B) = P(B)(1 - P(B)) = 0.5 \times 0.5 = 0.25
\end{equation}

Since $X_1 = 10 \cdot 1_A$ and $X_2 = 10 \cdot 1_B$:
\begin{align}
\text{Var}(X_1) &= 10^2 \times \text{Var}(1_A) = 100 \times 0.16 = 16\\
\text{Var}(X_2) &= 10^2 \times \text{Var}(1_B) = 100 \times 0.25 = 25
\end{align}

\subsection{Step 2: Find $\text{Cov}(X_1, X_2)$}

\begin{align}
\text{Cov}(X_1, X_2) &= \text{Cov}(10 \cdot 1_A, 10 \cdot 1_B)\\
&= 100 \times \text{Cov}(1_A, 1_B)
\end{align}

From the covariance formula:
\begin{align}
\text{Cov}(1_A, 1_B) &= E[1_A \cdot 1_B] - E[1_A]E[1_B]\\
&= P(A \cap B) - P(A)P(B)\\
&= 0.48 - 0.8 \times 0.5\\
&= 0.48 - 0.4\\
&= 0.08
\end{align}

Therefore:
\begin{equation}
\text{Cov}(X_1, X_2) = 100 \times 0.08 = 8
\end{equation}

\subsection{Step 3: Apply the Variance of Sum Formula}

\begin{tcolorbox}[colback=purple!5!white,colframe=purple!75!black,title=Variance of Sum Formula]
\begin{equation}
\text{Var}(X_1 + X_2) = \text{Var}(X_1) + \text{Var}(X_2) + 2\text{Cov}(X_1, X_2)
\end{equation}
\end{tcolorbox}

Substituting our values:
\begin{align}
\text{Var}(X) &= \text{Var}(X_1) + \text{Var}(X_2) + 2\text{Cov}(X_1, X_2)\\
&= 16 + 25 + 2(8)\\
&= 16 + 25 + 16\\
&= 57
\end{align}

\textbf{This matches our previous answer!} This method is often cleaner and more intuitive, especially when you need to work with sums of random variables.

\section{Key Takeaways}

\subsection{Why Not Independent?}

You might wonder: why can't we use $\text{Var}(X + Y) = \text{Var}(X) + \text{Var}(Y)$?

\textbf{Answer:} That formula only works when $X$ and $Y$ are \textbf{independent}. Here, $1_A$ and $1_B$ are \textbf{dependent} because:
\begin{equation}
P(B|A) = 0.6 \neq 0.5 = P(B)
\end{equation}

\subsection{General Formula for Variance of Sum}

For any two random variables $X$ and $Y$:
\begin{equation}
\text{Var}(X + Y) = \text{Var}(X) + \text{Var}(Y) + 2\text{Cov}(X, Y)
\end{equation}

where
\begin{equation}
\text{Cov}(X, Y) = E[XY] - E[X]E[Y]
\end{equation}

For our problem:
\begin{align}
\text{Cov}(1_A, 1_B) &= E[1_A \cdot 1_B] - E[1_A]E[1_B]\\
&= P(A \cap B) - P(A)P(B)\\
&= 0.48 - 0.8 \times 0.5\\
&= 0.48 - 0.4\\
&= 0.08 > 0
\end{align}

The positive covariance indicates that the two questions are positively correlated (doing well on one increases the chance of doing well on the other).

\subsection{Summary of Key Facts}

\begin{enumerate}
    \item \textbf{Indicator squared equals itself:} $1_A^2 = 1_A$

    \item \textbf{Product of indicators:} $1_A \cdot 1_B = 1_{A \cap B}$

    \item \textbf{Expectation of indicator:} $E[1_A] = P(A)$

    \item \textbf{Expectation of product:} $E[1_A \cdot 1_B] = P(A \cap B)$

    \item \textbf{Variance formula:} $\text{Var}(X) = E[X^2] - (E[X])^2$

    \item \textbf{Linearity of expectation:} Always works (no independence needed)

    \item \textbf{Variance of sum:} Only additive when variables are independent
\end{enumerate}

\end{document}
